\documentclass{report}
\usepackage{amsmath,amssymb}
\setlength{\parindent}{0mm}
\setlength{\parskip}{1em}
\begin{document}
\begin{center}
\rule{6in}{1pt} \
{\large Santiago Badia \\
{\bf Scalable space-time balancing domain-decomposition solvers}}

Universitat Polit\`ecnica de Catalunya (UPC) \\ Centre Internacional de M\`etodes Num\`erics a l'Enginyeria (CIMNE) \\ 08860 Castelldefels (Barcelona \\ Spain)
\\
{\tt sbadia@cimne.upc.edu}\\
Marc Olm\end{center}

The usual approach to transient problems is to exploit sequentiality in
time, and solve one space problem every time step. This approach has
recently been re-considered, motivated by the forthcoming exascale
supercomputers with billions of cores. This sequential approach has a
clear problem, parallelization cannot be exploited in time. Many key
computational engineering problems, e.g., turbulent flow simulations,
involve thousands to millions of time steps, and a scalable parallel
solver in space leads can lead to unacceptable computation times in these
problems. On the other hand, space parallelization always saturates at
some point and to efficiently exploit the forthcoming exascale platforms
one needs to exhibit further concurrency and less synchronization. In
this situation, the development of space-time solvers is an excellent
approach because deals with all time steps at once and exhibits much more
concurrency than space-only solvers. In this presentation, we will
introduce a new type of (non)linear space-time solvers for the solution
of finite element systems arising from the discretization of transient
problems, based on the novel extension of balancing domain decomposition
methods to space-time, and show numerically the scalability of the
proposed schemes on a set of academical problems, based on a highly
scalable overlapped multilevel task implementation.


\end{document}
